%!TEX root = main.tex

\subsection{Black-Rangarajan Duality}
\label{sec:BRduality}

This section revisits the Black-Rangarajan duality between robust estimation and outlier process~\cite{Black96ijcv-unification}. This theory is less known in robotics, and its applications have been mostly
targeting early vision problems, with few notable exceptions. 

% \marginpar{this  can be more formal}

\begin{lemma}[Black-Rangarajan Duality~\cite{Black96ijcv-unification}]
\label{lem:BRduality}
Given a robust cost function $\rho(\cdot)$, define $\phi(z) \doteq \rho(\sqrt{z})$. 
If $\phi(z)$ satisfies $\lim_{z\rightarrow 0} \phi'(z) = 1$, $\lim_{z\rightarrow \infty} \phi'(z) = 0$, and $\phi''(z) < 0$, then the robust estimation problem~\eqref{eq:generalEstimation}
is equivalent to % the following optimization % with outlier process:
%
\bea \label{eq:optiOutlierProcess} 
\min_{\vxx \in \calX, w_i \in [0,1]} \sumAllPointsi  \left[ w_i r^2(\vy_i,\vxx) + \Phi_\rho(w_i) \right],
\eea
%
where $ w_i \in [0,1]$ ($i=1,\dots,N$) are slack variables (or \emph{weights}) associated to each measurement $\vy_i$, 
and the 
%the weight (confidence level) for each measurement $\vy_i$. The 
  function $\Phi_\rho(w_i)$ (the so called \emph{outlier process}) defines a penalty on the weight $w_i$. 
  The  expression of $\Phi_\rho(w_i)$ depends on the choice of robust cost function $\rho(\cdot)$. 
\end{lemma}

The conditions on $\rho(\cdot)$ are satisfied by all common choices of robust costs~\cite{Black96ijcv-unification}.  
Besides presenting this fundamental result, asserting the equivalence between the outlier process~\eqref{eq:optiOutlierProcess}  and the robust estimation~\eqref{eq:generalEstimation}, 
Black and Rangarajan provide a procedure (see Fig.~10 in~\cite{Black96ijcv-unification}) to compute $\Phi_\rho(w_i)$ 
and show that common robust cost functions admit a simple analytical expression for $\Phi_\rho(w_i)$. 
% derive the expression 
% In addition, using the procedure in~\cite{Black96ijcv-unification}, $\Phi(w_i)$ can be derived in closed form for most robust functions $\rho(\cdot)$.
Interestingly, the outlier process~\eqref{eq:optiOutlierProcess} has been often used in robotics and SLAM~\cite{Agarwal13icra,Sunderhauf12iros}, without acknowledging the connection with robust estimation, and 
 with a heuristic design of the penalty terms $\Phi_\rho(w_i)$. 

Despite the elegance of~\prettyref{lem:BRduality}, Problem~\eqref{eq:optiOutlierProcess} remains hard to solve, 
due to its non-convexity. 
Approaches in robotics (\eg~\cite{Agarwal13icra,Sunderhauf12iros}) apply local optimization from an initial guess, 
resulting in brittle solutions (see~\cite{Carlone18ral-robustPGO2D} and Section~\ref{sec:applications}). 
\isExtended{To date, the only provably optimal solvers for~\eqref{eq:optiOutlierProcess} (when $\rho(\cdot)$ is a truncated least square cost) are given in~\cite{Yang19rss-teaser,Yang19iccv-QUASAR,Lajoie19ral-DCGM}.}{}
 %In the context of 3D perception,  
% However, although recently developed global solvers based on SDP and SOS relaxations are able to handle the non-convexity arising from feasible set $\calX$, they rely on least squares formulations ($w_i=1$ for all measurements) and fail in the presence of non-convexity in the cost function $\rho(\cdot)$. Therefore, solving the equivalent optimization~\eqref{eq:optiOutlierProcess} still requires a good initialization on either the weights $w_i$'s or the unknown variable $\vxx$~\cite{Black96ijcv-unification}.
 % allows compu
 % helps us handle the non-convexity in the cost function.

% Short section:
% \LC{
% \bit
% \item every robust cost can be converted to an outlier process, and for popular robust cost functions, the conversion is in closed form
% \item an intuitive idea is to alternate updates on variables
% \item unfortunately, the cost is nonconvex so this strategy fails (while it would succeed in a convex function)
% \eit
% }
